%!TEX root = ../template.tex %
% !TeX spellcheck = en_US
%%%%%%%%%%%%%%%%%%%%%%%%%%%%%%%%%%%%%%%%%%%%%%%%%%%%%%%%%%%%%%%%%%%%
%% chapter4.tex
%% NOVA thesis document file
%%
%% Chapter with lots of dummy text
%%%%%%%%%%%%%%%%%%%%%%%%%%%%%%%%%%%%%%%%%%%%%%%%%%%%%%%%%%%%%%%%%%%%

\typeout{NT FILE chapter4.tex}%

\chapter{Excitonic effects in Two-dimensional Semiconductors}
\label{ch:Excitons}

\section{Introduction}
\label{sec:introduction}

Here I give an explanation of the excitons from a phenomenological point of view, and Wannier equation, the simplest model.

\section{The Bethe--Salpeter Equation: application to effective models}
\label{sec:bethe_salpeter_equation_effective_models}

Here I introduce Nuno's approach to excitons and the Bethe--Salpeter equation.
We need this since we used it to compute the renormalized band gap in the ACS Nano.
We can also derive back again the Wannier equation.

\section{The Bethe--Salpeter Equation: exact diagonalization}
\label{sec:bethe_salpeter_equation_exact}

Here I discuss the excitonic problem using numerical methods of the xatu paper, emphasizing the TB approximation in the methods.
First with Rytova-Keldysh potential, already introduced before.
Then with the microscopic dielectric function introduced previously, to be computed later.

\section{Optical effects in 2D semiconductors}
\label{sec:optical_effects}

Here we provide the theory for the optical conductivity.

\subsection{Non-interacting optical conductivity}
\label{subsec:non_interacting_conductivity}

Here I could do a short derivation, that is in the book by G.Grosso and G.Parravicini.

\subsection{Excitonic optical conductivity}
\label{subsec:excitonic_conductivity}

Here I have to look at the thesis of Joao Henriques, previous student of Nuno.